\section{Метод FieldMathCalculations::Calc\_Circulation\_Stokes}
\footnotesize
\begin{listing}{1}
/*************************************************************************
* Метод         : FieldMathCalculations::Calc_Circulation_Stokes         *
*------------------------------------------------------------------------*
* Назначение    : Расчет циркуляции векторного поля через теорему Стокса *
* Входные       : Нет                                                    *
* Выходные      : double - значение циркуляции                           *
*------------------------------------------------------------------------*
* ОПИСАНИЕ ФУНКЦИОНИРОВАНИЯ:                                             *
* Вычисляет циркуляцию поля по контуру, используя теорему Стокса (переход*
* от криволинейного интеграла к интегралу от ротора по поверхности).     *
* *
* АЛГОРИТМ:                                                              *
* 1. Интегрирование ротора поля по поверхности через IntegrateQuarterDisk. *
* 2. Вычисление проекции вектора ротора на нормаль к поверхности.         *
*************************************************************************/
double FieldMathCalculations::Calc_Circulation_Stokes() {
    return IntegrateQuarterDisk([](double x, double y) {
        double z = sqrt(fmax(0.0, 1.0 - x * x - y * y));
        if (z < 1e-5) return 0.0;
        Vector3D rot = Calc_Rotor(EvaluateFieldA, x, y, z);
        return (rot.x * x + rot.y * y + rot.z * z) / z;
    });
}
\end{listing}
\normalsize

\pagebreak

\section{Метод FieldMathCalculations::Calc\_Circulation\_Direct}
\footnotesize
\begin{listing}{1}
/*************************************************************************
* Метод         : FieldMathCalculations::Calc_Circulation_Direct         *
*------------------------------------------------------------------------*
* Назначение    : Прямое вычисление циркуляции                          *
* Входные       : Нет                                                    *
* Выходные      : double - значение циркуляции                           *
*------------------------------------------------------------------------*
* ОПИСАНИЕ ФУНКЦИОНИРОВАНИЯ:                                             *
* Выполняет непосредственное интегрирование поля вдоль граничных         *
* отрезков в первой октанте (три стороны контура).                       *
* *
* АЛГОРИТМ:                                                              *
* 1. Параметризация трех сторон контура.                                 *
* 2. Численное суммирование скалярного произведения поля на касательный  *
* вектор (dl) с шагом dphi.                                              *
*************************************************************************/
double FieldMathCalculations::Calc_Circulation_Direct() {
    int steps = 100; 
    double dphi = (M_PI / 2.0) / steps, sum = 0.0;
    for (int i = 0; i < steps; i++) {
        double phi = (i + 0.5) * dphi;
        Vector3D a1 = EvaluateFieldA(cos(phi), sin(phi), 0.0);
        sum += (a1.x * (-sin(phi)) + a1.y * cos(phi));
        Vector3D a2 = EvaluateFieldA(0.0, cos(phi), sin(phi));
        sum += (a2.y * (-sin(phi)) + a2.z * cos(phi));
        Vector3D a3 = EvaluateFieldA(sin(phi), 0.0, cos(phi));
        sum += (a3.z * (-sin(phi)) + a3.x * cos(phi));
    }
    return sum * dphi;
}
\end{listing}
\normalsize

\pagebreak

\section{Метод FieldMathCalculations::GetAnalyticalPotential}
\footnotesize
\begin{listing}{1}
/*************************************************************************
* Метод         : FieldMathCalculations::GetAnalyticalPotential          *
*------------------------------------------------------------------------*
* Назначение    : Аналитическое нахождение потенциала поля B             *
* Входные       : Нет                                                    *
* Выходные      : Polynomial - результирующая функция потенциала         *
*------------------------------------------------------------------------*
* ОПИСАНИЕ ФУНКЦИОНИРОВАНИЯ:                                             *
* Вычисляет скалярный потенциал U(x,y,z), такой что grad(U) = B.         *
* *
* АЛГОРИТМ:                                                              *
* 1. Интегрирование компонент (Bx по x, By по y, Bz по z).               *
* 2. Использование метода SetVariableToZero для исключения перекрестных  *
* членов, чтобы избежать дублирования данных при восстановлении.        *
* 3. Алгебраическое сложение полученных функций (Add).                   *
*************************************************************************/
Polynomial FieldMathCalculations::GetAnalyticalPotential() {
    Polynomial U1 = GetFieldBX().SetVariableToZero(false, true, true)
                                .IntegrateSymbolically(0);
    Polynomial U2 = GetFieldBY().SetVariableToZero(false, false, true)
                                .IntegrateSymbolically(1);
    Polynomial U3 = GetFieldBZ().SetVariableToZero(false, false, false)
                                .IntegrateSymbolically(2);
    Polynomial U; U.Add(U1); U.Add(U2); U.Add(U3);
    return U;
}
\end{listing}
\normalsize

\pagebreak

\section{Метод FieldMathCalculations::VerifyPotentialCorrectness}
\footnotesize
\begin{listing}{1}
/*************************************************************************
* Метод         : FieldMathCalculations::VerifyPotentialCorrectness      *
*------------------------------------------------------------------------*
* Назначение    : Верификация точности найденного потенциала             *
* Входные       : Нет                                                    *
* Выходные      : bool - статус корректности                             *
*------------------------------------------------------------------------*
* ОПИСАНИЕ ФУНКЦИОНИРОВАНИЯ:                                             *
* Проверяет условие grad(U) = B в контрольных точках сетки.              *
* *
* АЛГОРИТМ:                                                              *
* 1. В контрольных точках вычисляются частные производные от U.          *
* 2. Сравнение градиента с вектором поля B. Если разница превышает       *
* порог (0.05) в любой точке, возвращается false.                       *
*************************************************************************/
bool FieldMathCalculations::VerifyPotentialCorrectness() {
    Polynomial U = GetAnalyticalPotential();
    FUNC3D U_eval = [&](double px, double py, double pz) { 
        return U.Evaluate(px, py, pz); 
    };
    for (double x = 0.2; x <= 0.8; x += 0.3) {
        for (double y = 0.2; y <= 0.8; y += 0.3) {
            double du_dx = Calc_PartialDerivative(U_eval, 1, x, y, 0.5);
            double du_dy = Calc_PartialDerivative(U_eval, 2, x, y, 0.5);
            Vector3D valB = EvaluateFieldB(x, y, 0.5);
            if (abs(du_dx - valB.x) > 0.05 || abs(du_dy - valB.y) > 0.05) 
                return false;
        }
    }
    return true;
}
\end{listing}
\normalsize

\pagebreak

\section{Метод FieldMathCalculations::CheckRotorBZero}
\footnotesize
\begin{listing}{1}
/*************************************************************************
* Метод         : FieldMathCalculations::CheckRotorBZero                 *
*------------------------------------------------------------------------*
* Назначение    : Проверка потенциальности поля B (rot B = 0)            *
* Входные       : Нет                                                    *
* Выходные      : bool - результат проверки                              *
*------------------------------------------------------------------------*
* ОПИСАНИЕ ФУНКЦИОНИРОВАНИЯ:                                             *
* Численная проверка условия отсутствия вихрей в поле B.                 *
* *
* АЛГОРИТМ:                                                              *
* 1. Сканирование области по сетке (x, y).                               *
* 2. Вычисление ротора в каждой точке.                                   *
* 3. Если модуль ротора превышает 1e-3, поле считается непотенциальным.  *
*************************************************************************/
bool FieldMathCalculations::CheckRotorBZero() {
    for (double x = 0.2; x <= 0.8; x += 0.3) {
        for (double y = 0.2; y <= 0.8; y += 0.3) {
            Vector3D r = Calc_Rotor(EvaluateFieldB, x, y, 0.5);
            if (abs(r.x) > 1e-3 || abs(r.y) > 1e-3 || abs(r.z) > 1e-3) 
                return false;
        }
    }
    return true;
}
\end{listing}
\normalsize

\pagebreak
